\documentclass[10pt,a4paper]{amsart}
\usepackage{amsmath,amssymb,amsthm,graphicx}
\usepackage{verbatim}
\newenvironment{reference}{\comment}{\endcomment}
%\newenvironment{reference}{}{}
\newenvironment{slogan}{\comment}{\endcomment}
\newenvironment{history}{\comment}{\endcomment}

% For commutative diagrams we use Xy-pic
\usepackage[all]{xy}

% We use 2cell for 2-commutative diagrams.
\xyoption{2cell}
% Theorem environments.
%
\theoremstyle{plain}
\newtheorem{theorem}[subsection]{Theorem}
\newtheorem{proposition}[subsection]{Proposition}
\newtheorem{lemma}[subsection]{Lemma}

\theoremstyle{definition}
\newtheorem{definition}[subsection]{Definition}
\newtheorem{example}[subsection]{Example}
\newtheorem{exercise}[subsection]{Exercise}
\newtheorem{situation}[subsection]{Situation}

\theoremstyle{remark}
\newtheorem{remark}[subsection]{Remark}
\newtheorem{remarks}[subsection]{Remarks}

\numberwithin{equation}{subsection}

% Macros
%
\def\lim{\mathop{\mathrm{lim}}\nolimits}
\def\colim{\mathop{\mathrm{colim}}\nolimits}
\def\Spec{\mathop{\mathrm{Spec}}}
\def\Hom{\mathop{\mathrm{Hom}}\nolimits}
\def\Ext{\mathop{\mathrm{Ext}}\nolimits}
\def\SheafHom{\mathop{\mathcal{H}\!\mathit{om}}\nolimits}
\def\SheafExt{\mathop{\mathcal{E}\!\mathit{xt}}\nolimits}
\def\Sch{\mathit{Sch}}
\def\Mor{\mathop{\mathrm{Mor}}\nolimits}
\def\Ob{\mathop{\mathrm{Ob}}\nolimits}
\def\Sh{\mathop{\mathit{Sh}}\nolimits}
\def\NL{\mathop{N\!L}\nolimits}
\def\CH{\mathop{\mathrm{CH}}\nolimits}
\def\proetale{{pro\text{-}\acute{e}tale}}
\def\etale{{\acute{e}tale}}
\def\QCoh{\mathit{QCoh}}
\def\Ker{\mathop{\mathrm{Ker}}}
\def\Im{\mathop{\mathrm{Im}}}
\def\Coker{\mathop{\mathrm{Coker}}}
\def\Coim{\mathop{\mathrm{Coim}}}

% Boxtimes
%
\DeclareMathSymbol{\boxtimes}{\mathbin}{AMSa}{"02}

%
% Macros for moduli stacks/spaces
%
\def\QCohstack{\mathcal{QC}\!\mathit{oh}}
\def\Cohstack{\mathcal{C}\!\mathit{oh}}
\def\Spacesstack{\mathcal{S}\!\mathit{paces}}
\def\Quotfunctor{\mathrm{Quot}}
\def\Hilbfunctor{\mathrm{Hilb}}
\def\Curvesstack{\mathcal{C}\!\mathit{urves}}
\def\Polarizedstack{\mathcal{P}\!\mathit{olarized}}
\def\Complexesstack{\mathcal{C}\!\mathit{omplexes}}
% \Pic is the operator that assigns to X its picard group, usage \Pic(X)
% \Picardstack_{X/B} denotes the Picard stack of X over B
% \Picardfunctor_{X/B} denotes the Picard functor of X over B
\def\Pic{\mathop{\mathrm{Pic}}\nolimits}
\def\Picardstack{\mathcal{P}\!\mathit{ic}}
\def\Picardfunctor{\mathrm{Pic}}
\def\Deformationcategory{\mathcal{D}\!\mathit{ef}}
\usepackage{cleveref}
\crefname{originalreference}{original reference}{original references}
\Crefname{originalreference}{Original reference}{Original references}
\makeatletter
\newlabel{lemma-derived-quasi-coherent-small-etale-site}{{071Q}{source}{}{sourcepointer-1}{}}
\newlabel{lemma-derived-quasi-coherent-small-etale-site@cref}{{[originalreference][0][]071Q}{[1][0][]source}{}{}{}}
\makeatother
\crefname{originalreference}{original reference}{original references}
\Crefname{originalreference}{Original reference}{Original references}
\makeatletter
\newlabel{lemma-descend-perfect}{{08HG}{source}{}{sourcepointer-2}{}}
\newlabel{lemma-descend-perfect@cref}{{[originalreference][0][]08HG}{[1][0][]source}{}{}{}}
\makeatother
\crefname{originalreference}{original reference}{original references}
\Crefname{originalreference}{Original reference}{Original references}
\makeatletter
\newlabel{lemma-glue}{{01JB}{source}{}{sourcepointer-3}{}}
\newlabel{lemma-glue@cref}{{[originalreference][0][]01JB}{[1][0][]source}{}{}{}}
\makeatother
\crefname{originalreference}{original reference}{original references}
\Crefname{originalreference}{Original reference}{Original references}
\makeatletter
\newlabel{lemma-induction-principle-enlarge}{{09IT}{source}{}{sourcepointer-4}{}}
\newlabel{lemma-induction-principle-enlarge@cref}{{[originalreference][0][]09IT}{[1][0][]source}{}{}{}}
\makeatother
\crefname{originallemma}{original lemma}{original lemmas}
\Crefname{originallemma}{Original lemma}{Original lemmas}
\makeatletter
\newlabel{perfect-lemma-lift-map}{{08EH}{source}{}{sourcepointer-5}{}}
\newlabel{perfect-lemma-lift-map@cref}{{[originallemma][0][]08EH}{[1][0][]source}{}{}{}}
\makeatother
\crefname{originallemma}{original lemma}{original lemmas}
\Crefname{originallemma}{Original lemma}{Original lemmas}
\makeatletter
\newlabel{perfect-lemma-lift-perfect-complex-plus-shift-support}{{08EK}{source}{}{sourcepointer-6}{}}
\newlabel{perfect-lemma-lift-perfect-complex-plus-shift-support@cref}{{[originallemma][0][]08EK}{[1][0][]source}{}{}{}}
\makeatother
\crefname{originalreference}{original reference}{original references}
\Crefname{originalreference}{Original reference}{Original references}
\makeatletter
\newlabel{remark-match-total-direct-images}{{08GH}{source}{}{sourcepointer-7}{}}
\newlabel{remark-match-total-direct-images@cref}{{[originalreference][0][]08GH}{[1][0][]source}{}{}{}}
\makeatother
\title{Lifting perfect maps supported on T from a quasi-compact open to X}
\author{The Stacks Project Authors}
\date{stacks-project:spaces-perfect.tex:local-snapshot-20261012; GNU FDL 1.2 or later}
\begin{document}
\maketitle
\noindent\textit{Editorial note.} All delivered mathematical text is the author's own, copied from the pinned Stacks Project chapter file spaces-perfect.tex; no statement or proof step was rewritten, reordered or removed.\par\medskip
\noindent\textit{Editorial note.} The counted proof is the author's complete argument: the proof reduces to an elementary distinguished square by the induction principle, lifts the perfect object and its map over the affine piece using the two scheme lemmas, and gluees the resulting morphism via lemma-glue\par\medskip
\noindent\textit{Editorial note.} The label lemma-lift-map-from-perfect-complex-with-support names the counted statement, and the proof environment immediately following it in the source is delivered verbatim as its proof.\par\medskip
\noindent\textit{Editorial note.} The counted statement is carried by the author's own lemma environment in the chapter source.\par\medskip
\noindent\textit{Allowed standard prerequisites (editorial).}\begin{itemize}
\item the induction principle for enlarging the open part (lemma-induction-principle-enlarge of the same chapter) reducing to elementary distinguished squares
\item the comparison lemma between the derived category of the small etale site and quasi-coherent modules (lemma-derived-quasi-coherent-small-etale-site of the same chapter) with the compatibilities of remark-match-total-direct-images
\item perfectness agrees between the Zariski and etale sites (lemma-descend-perfect of the same chapter)
\item the scheme-case lifting of perfect complexes with support over affine covers (Derived Categories of Schemes, Lemma perfect-lemma-lift-perfect-complex-plus-shift-support)
\item the lifting of maps along perfect complexes after tensoring with an injective resolution (Derived Categories of Schemes, Lemma perfect-lemma-lift-map)
\item the glueing of morphisms in D(O\_X) from W and V over the intersection (lemma-glue of the same chapter)
\end{itemize}\medskip
\ifdefined\hypertarget\hypertarget{sourcepointer-1}{}\fi
\noindent\textit{Source reference.} Original reference 071Q: sources/stacks-spaces-perfect/pinned-original-source.tex, source lines 402--402. Prerequisite pointer invoked at proof line(s) 3834; it names a standard fact of the same chapter that the counted argument uses, and every step of the counted proof built on it is retained verbatim in the delivered body.\par\medskip
\ifdefined\hypertarget\hypertarget{sourcepointer-2}{}\fi
\noindent\textit{Source reference.} Original reference 08HG: sources/stacks-spaces-perfect/pinned-original-source.tex, source lines 3288--3288. Prerequisite pointer invoked at proof line(s) 3838; it names a standard fact of the same chapter that the counted argument uses, and every step of the counted proof built on it is retained verbatim in the delivered body.\par\medskip
\ifdefined\hypertarget\hypertarget{sourcepointer-3}{}\fi
\noindent\textit{Source reference.} Original reference 01JB: sources/stacks-spaces-perfect/pinned-original-source.tex, source lines 2206--2206. Prerequisite pointer invoked at proof line(s) 3855; it names a standard fact of the same chapter that the counted argument uses, and every step of the counted proof built on it is retained verbatim in the delivered body.\par\medskip
\ifdefined\hypertarget\hypertarget{sourcepointer-4}{}\fi
\noindent\textit{Source reference.} Original reference 09IT: sources/stacks-spaces-perfect/pinned-original-source.tex, source lines 1707--1707. Prerequisite pointer invoked at proof line(s) 3829; it names a standard fact of the same chapter that the counted argument uses, and every step of the counted proof built on it is retained verbatim in the delivered body.\par\medskip
\ifdefined\hypertarget\hypertarget{sourcepointer-5}{}\fi
\noindent\textit{Source reference.} Original lemma 08EH: sources/stacks-chow/tags.txt, source lines 10673--10673. Prerequisite pointer invoked at proof line(s) 3846; it names a standard fact of a sibling chapter that the counted argument uses, and every step of the counted proof built on it is retained verbatim in the delivered body.\par\medskip
\ifdefined\hypertarget\hypertarget{sourcepointer-6}{}\fi
\noindent\textit{Source reference.} Original lemma 08EK: sources/stacks-chow/tags.txt, source lines 10676--10676. Prerequisite pointer invoked at proof line(s) 3842; it names a standard fact of a sibling chapter that the counted argument uses, and every step of the counted proof built on it is retained verbatim in the delivered body.\par\medskip
\ifdefined\hypertarget\hypertarget{sourcepointer-7}{}\fi
\noindent\textit{Source reference.} Original reference 08GH: sources/stacks-spaces-perfect/pinned-original-source.tex, source lines 896--896. Prerequisite pointer invoked at proof line(s) 3835; it names a standard fact of the same chapter that the counted argument uses, and every step of the counted proof built on it is retained verbatim in the delivered body.\par\medskip
\begin{lemma}
\label{lemma-lift-map-from-perfect-complex-with-support}
Let $S$ be a scheme. Let $X$ be a quasi-compact and quasi-separated
algebraic space over $S$. Let $W \subset X$ be a quasi-compact open.
Let $T \subset |X|$ be a closed subset such that
$X \setminus T \to X$ is a quasi-compact morphism.
Let $E$ be an object of $D_\QCoh(\mathcal{O}_X)$.
Let $\alpha : P \to E|_W$ be a map where $P$ is a perfect object of
$D(\mathcal{O}_W)$ supported on $T \cap W$. Then there exists a map
$\beta : R \to E$ where $R$ is a perfect object of $D(\mathcal{O}_X)$
supported on $T$ such that $P$ is a direct summand of $R|_W$ in
$D(\mathcal{O}_W)$ compatible $\alpha$ and $\beta|_W$.
\end{lemma}

\begin{proof}
We will use the induction principle of
Lemma \ref{lemma-induction-principle-enlarge} to prove this.
Thus we immediately reduce to the case where we have an
elementary distinguished square $(W \subset X, f : V \to X)$
with $V$ affine and $P \to E|_W$ as in the statement of the lemma.
In the rest of the proof we will use
Lemma \ref{lemma-derived-quasi-coherent-small-etale-site}
(and the compatibilities of Remark \ref{remark-match-total-direct-images})
for the representable algebraic spaces $V$ and $W \times_X V$.
We will also use the fact that perfectness on the Zariski site
and \'etale site agree, see Lemma \ref{lemma-descend-perfect}.

\medskip\noindent
By Derived Categories of Schemes,
Lemma \ref{perfect-lemma-lift-perfect-complex-plus-shift-support}
we can choose a perfect object $Q$ in $D(\mathcal{O}_V)$
supported on $f^{-1}T$ and an isomorphism
$Q|_{W \times_X V} \to (P \oplus P[1])|_{W \times_X V}$. By
Derived Categories of Schemes, Lemma \ref{perfect-lemma-lift-map}
we can replace $Q$ by $Q \otimes^\mathbf{L} I$ (still supported on $f^{-1}T$)
and assume that the map
$$
Q|_{W \times_X V} \to (P \oplus P[1])|_{W \times V}
\longrightarrow P|_{W \times_X V}
\longrightarrow
E|_{W \times_X V}
$$
lifts to $Q \to E|_V$. By Lemma \ref{lemma-glue}
we find an morphism $a : R \to E$ of $D(\mathcal{O}_X)$
such that $a|_W$ is isomorphic to $P \oplus P[1] \to E|_W$
and $a|_V$ isomorphic to $Q \to E|_V$.
Thus $R$ is perfect and supported on $T$ as desired.
\end{proof}

\end{document}
